problem:  
Let \((M^{2n+1}, T^{1,0}M, \theta)\) be a compact, strictly pseudoconvex, **pseudo‑Einstein** pseudohermitian manifold with CR dimension \(n\ge2\).  
Under a conformal change \(\tilde{\theta}=e^{2u}\theta\), denote by \(\tilde{\Delta}_b\) the sub‑Laplacian and by \(\lambda_1(\tilde{\Delta}_b)\) its first nonzero eigenvalue.  
On the standard pseudohermitian sphere \((S^{2n+1},\theta_{\mathrm{std}})\), equality holds in the CR Lichnerowicz–Obata bound:  
\[
\lambda_1(\Delta_b) = \tfrac{n}{2n+2}\,R_{\theta_{\mathrm{std}}}.
\]

**Problem (Quantitative CR Obata‑type rigidity):**  
Fix a normalization \(\mathrm{Vol}(M,\tilde{\theta})=1\) for all conformal changes.  
Suppose there exists a sequence \(\tilde{\theta}_k=e^{2u_k}\theta\) such that  

1. the Tanaka–Webster torsion satisfies \( \|A_{\alpha\beta}(\tilde{\theta}_k)\|_{L^\infty}\le C\);  
2. the Webster scalar curvature \(\tilde{R}_k>0\);  
3. the spectral ratio approaches the spherical constant:
   \[
   \frac{\lambda_1(\tilde{\Delta}_{b,k})}{\min_M \tilde{R}_k}\longrightarrow \frac{n}{2n+2}.
   \]  

Does it follow, after passing to a subsequence and scaling, that \((M,\theta)\) is CR equivalent to the standard pseudohermitian sphere \((S^{2n+1},\theta_{\mathrm{std}})\)?  

Equivalently, can a non‑spherical, pseudo‑Einstein CR manifold admit a **spectrally nearly extremal** sequence of conformal pseudohermitian structures with bounded torsion and unit volume?

---

Why it is a "good" problem:

1. **Bridges spectral analysis and CR geometry.**  
   The problem asks whether near saturation of the spectral bound for the sub‑Laplacian forces global geometric rigidity. It intertwines subelliptic spectral theory and global CR structure, potentially uniting analytical methods (eigenvalue estimates, blow‑up analysis, compactness theorems) with pseudohermitian geometry.

2. **Extension of classical rigidity into a quantitative regime.**  
   The Jerison–Lee CR Obata theorem covers the exact equality case; this question addresses the stability of that equality. Establishing a quantitative version parallels key developments in Riemannian geometry but for subelliptic operators, creating a new research direction in **quantitative pseudohermitian rigidity**.

3. **Analytic depth with novel geometric constraints.**  
   Controlling curvature, torsion, and conformal factors simultaneously introduces complex nonlinear interactions. Determining whether such asymptotic extremals can exist will require new compactness or energy‑concentration principles in the CR setting. The pseudo‑Einstein assumption provides geometric coherence for curvature transformations, while bounded torsion prevents degenerate limits.

4. **Simple but profound formulation.**  
   The question “Does near equality in the CR Lichnerowicz–Obata bound force the sphere?” is elementary to state yet engages deep analytic and geometric mechanisms. A positive answer would yield a new stability theorem; a counterexample would uncover unexpected flexibility in CR manifolds—either outcome propelling understanding of the geometric–analytic boundary within pseudohermitian geometry.